% Part: lambda-calculus
% Chapter: lambda-definability
% Section: pairs

\documentclass[../../../include/open-logic-section]{subfiles}

\begin{document}

\olfileid{lam}{ldf}{pai}
\olsection{序对与前驱}

\begin{defn}
$M$ 与 $N$ 的序对（记作 $\tuple{M,N}$）定义如下：
\[
\tuple{M,N} \ident \lambd[f][fMN].
\]
\end{defn}
  
直观上，它是一个接受函数作为输入，并将该函数作用于序对的两个元素的函数。
顺着这个思路，我们便有了如下的构造子，它接受两个项并返回包含它们的序对：
\[
  \fn{Pair} \ident \lambd[mn][\lambd[f][fmn]]
\]
给定一个序对，我们也希望能提取其中的元素。
为此我们需要两个访问函数，它们接受一个序对作为自变元，并返回其中的第一个或第二个元素：
\begin{align*}
  \fn{Fst} & \ident \lambd[p][p(\lambd[mn][m])]\\
  \fn{Snd} & \ident \lambd[p][p(\lambd[mn][n])]
\end{align*}

\begin{prob}
  解释为什么访问函数 $\fn{Fst}$ 和 $\fn{Snd}$ 是有效的。
\end{prob}

现在，有了序对我们便可以 $\lambd$-定义前驱函数：
\[
  \fn{Pred} \ident \lambd[n][\fn{Fst}(n (\lambd[p][\tuple{\fn{Snd}\, {p}, \fn{Succ}(\fn{Snd}\, {p})}]) \tuple{\num 0, \num 0})]
\]
回想一下，$\num n\, f x$ 归约到 $f^{n}(x)$；
在这里，$f$ 是一个接受序对 $p$ 作为自变元并返回一个新序对的函数，新序对包含 $p$ 的第二个分量以及该分量的后继；
$x$ 是序对 $\tuple{0,0}$。
因此，当 $n=0$ 时结果为 $\tuple{0,0}$，否则为 $\tuple{\num{n-1}, \num n}$。
接着，$\fn{Pred}$ 返回结果的第一个分量。

减法可以定义为将 $\fn{Pred}$ 对 $a$ 作用 $b$ 次：
\[
\fn{Sub} \ident \lambd[ab][b \fn{Pred}\, a].
\]
    
\end{document}